\section{Dinamika Pasar Finansial}

\subsection{Sejarah Uang}
Manusia memiliki kebutuhan yang sangat beragam untuk menopang eksistensinya, namun secara biologis pemenuhan kebutuhan tersebut senantiasa memiliki batas kecukupan yang bersifat fundamental dalam ruang lingkup konsumsi harian. Berbeda dengan kebutuhan primer, keinginan manusia cenderung bersifat tidak terbatas (\textit{infinite}) sementara ketersediaan sumber daya alam yang dapat digunakan untuk memenuhinya bersifat sangat terbatas (\textit{scarcity}). Dalam konteks ini, uang hadir bukan sekadar sebagai alat transaksi, melainkan sebagai instrumen filosofis untuk membatasi keinginan manusia sekaligus menjadi penyeimbang terhadap ketersediaan sumber daya di alam. Evolusi sistem ekonomi dimulai dari mekanisme barter yang memiliki tingkat inefisiensi masif karena adanya kendala dalam pencocokan kebutuhan ganda antar individu yang terlibat dalam pertukaran. Penggunaan uang komoditas (\textit{commodity money}) seperti rempah-rempah atau kerang kemudian muncul sebagai solusi awal untuk menjembatani transaksi tersebut, meskipun efektivitasnya terhambat oleh masalah portabilitas fisik. Logam mulia kemudian diadopsi secara luas untuk menstandarisasi nilai tukar karena sifatnya yang langka dan tahan lama, meskipun penggunaannya masih menyisakan hambatan logistik yang berat dalam mobilisasi kapital. Transisi menuju uang \textit{fiat} menandai pergeseran paradigma di mana nilai sebuah mata uang tidak lagi bersandar pada fisik komoditas, melainkan pada janji bayar dan legitimasi otoritas pemerintah. Menurut \cite{kolmogorov_three_1968}, uang secara teoretis berfungsi sebagai alat kompresi data (\textit{data compression}) yang menyederhanakan rasio harga barter yang sangat kompleks menjadi satu standar nilai yang efisien secara informasional. Figure 1 secara komprehensif mengilustrasikan proses evolusi medium pertukaran dari sistem barter konvensional menuju representasi nilai digital yang terenkripsi secara kriptografis dalam jaringan global.

Pada fase awal perkembangannya, uang kertas atau uang \textit{fiat} digunakan secara terbatas sekadar sebagai surat bukti kepemilikan atas cadangan emas yang disimpan secara fisik di dalam brankas perbankan. Namun, seiring berjalannya waktu, institusi perbankan mulai mencetak lembaran kertas secara berlebihan hingga melampaui jumlah cadangan emas riil yang tersedia untuk menjamin nilai lembaran tersebut. Praktik ekspansi moneter yang tidak terkontrol ini mengakibatkan runtuhnya janji penukaran sistemik, di mana satu lembar kertas tidak lagi memiliki kepastian untuk ditukar dengan jumlah emas yang setara. Fenomena kegagalan standar emas tersebut memaksa paradigma teknologi moneter global untuk bergeser secara radikal dari sandaran nilai fisik menuju sistem nilai fungsional murni atau mata uang mengambang (\textit{floating currency}). Era modern kemudian melengkapi proses dematerialisasi ini melalui adopsi uang digital yang secara total menghapus batasan geografis serta wujud fisik dalam setiap aliran transaksi modal global. Perkembangan lebih lanjut mengenai uang kripto (\textit{cryptocurrency}) diproyeksikan akan menjadi standar mata uang di masa depan yang mengandalkan algoritma kriptografi untuk menjamin integritas dan kelangkaan nilai dalam ekosistem digital.

\subsection{Inflasi}
Penting untuk ditegaskan secara fundamental bahwa uang dalam sistem ekonomi modern berfungsi sebagai alat tukar (\textit{medium of exchange}) yang dinamis, bukan sekadar instrumen penyimpan kekayaan yang bersifat statis. Inflasi secara teknis didefinisikan sebagai fenomena kenaikan harga barang dan jasa secara umum yang mengakibatkan degradasi daya beli mata uang secara berkelanjutan dalam jangka waktu tertentu. Secara makroekonomi, fenomena ini mencerminkan adanya luapan likuiditas atau jumlah uang beredar yang jauh melampaui kapasitas ketersediaan barang dan jasa yang ditawarkan di pasar. Figure 2 menyajikan kurva penawaran dan permintaan yang mengilustrasikan pergeseran titik keseimbangan harga akibat tekanan inflasi yang timbul dari ketidakseimbangan sirkulasi likuiditas pasar. Volatilitas inflasi menjadi parameter risiko yang sangat krusial dalam analisis finansial karena sifatnya yang stokastik serta memiliki tingkat kesulitan yang sangat tinggi untuk diprediksi secara deterministik. Menurut \cite{engle_1982}, varians dari tingkat inflasi menunjukkan pola autokoreksi heteroskedastik yang sangat dipengaruhi oleh guncangan ekonomi masa lalu, sehingga menciptakan ketidakpastian dalam estimasi nilai di masa depan.

Dampak dari fenomena inflasi merambat secara luas ke sektor investasi melalui distorsi sinyal harga serta pergeseran ekspektasi terhadap tingkat imbal hasil riil yang akan diterima oleh investor. Kekayaan riil seorang agen ekonomi secara otomatis akan mengalami pengikisan apabila tingkat inflasi tahunan berada di atas tingkat suku bunga tabungan atau imbal hasil dari instrumen kas lainnya. Dalam perspektif ini, inflasi bertindak sebagai "pajak tersembunyi" (\textit{hidden tax}) yang secara sistemik memaksa para pelaku pasar untuk mencari instrumen investasi produktif yang mampu menghasilkan imbal hasil di atas laju kenaikan harga. Dinamika ini menjadi landasan bagi bank sentral dalam merumuskan kebijakan moneter yang bertujuan untuk menjaga stabilitas sirkulasi uang serta memitigasi risiko devaluasi daya beli masyarakat. Oleh karena itu, diperlukan sebuah sistem manajemen kekayaan yang lebih realistis dan adaptif guna menjaga nilai kapital agar tetap relevan di tengah perubahan kondisi ekonomi yang bersifat fluktuatif.
